% ============================================================
% Capítulo: Conclusiones (Bloque B.14)
% ============================================================
\chapter{Conclusiones}
\label{cap:conclusiones}

\section{Respuesta condensada a las preguntas de investigación}
\label{sec:concl-rqs}

\textbf{RQ1 (rendimiento).} El sistema cumple ambos umbrales de latencia
exigidos con margen amplio (p95 $19{,}49$\,ms caliente,
$7{,}50$\,ms frío), aunque la comparación entre escenarios no aísla
limpiamente el efecto del caché por una limitación metodológica ya
declarada -- el umbral se cumple, la atribución causal exacta queda
matizada.

\textbf{RQ2 (seguridad).} De 6 controles OWASP auditados, 4 presentaron
un hallazgo real corregido dentro del mismo ciclo ($66{,}7$\,\%), y 2 se
verificaron sin hallazgo real de código (A03 desde la corrida original;
A02 tras completarse la verificación de TLS/HSTS contra el despliegue
público real, 18-ago-2026) -- los 6 controles quedan con resultado
cerrado; la auditoría encontró y forzó a corregir defectos concretos, no
fue un ejercicio cosmético.

\textbf{RQ3 (usabilidad).} Sin respuesta todavía: $N=0$ de los
$N\geq15$ participantes exigidos por la guía. Queda como trabajo en
progreso, bloqueado en la cadena de dependencias correcta (despliegue
público $\to$ reclutamiento $\to$ aplicación del instrumento SUS), no
como una omisión.

\textbf{RQ4 (arquitectura híbrida).} La estrategia CRUD-ORM/SP no
perjudicó la trazabilidad ($48{,}8$\,\% CRUD-ORM, $18{,}6$\,\% SP,
$4{,}7$\,\% híbrido, distribución medida sobre los 43 requisitos
originales de \autoref{sec:matriz-resumen} -- la matriz creció a 46
requisitos tras el cierre, ver \autoref{cap:ingenieria-requisitos}; el
desglose porcentual no se recalculó para los 3 nuevos por alcance de
esta corrección) ni la cobertura de pruebas ($82{,}97$\,\% de líneas),
pero tuvo un costo real de implementación: una migración forzada de
\texttt{@Procedure} a \texttt{@Query} nativa por un defecto documentado
de Hibernate con procedimientos multi-\texttt{OUT}.

\section{Contribuciones -- recapituladas con evidencia del capítulo de resultados}
\label{sec:concl-contribuciones}

El \autoref{cap:introduccion} enumeró cinco contribuciones de este
proyecto. Con la evidencia empírica del \autoref{cap:resultados} ya
reunida, se sostienen así:

\begin{itemize}[leftmargin=1.2cm]
  \item Un sistema web funcional con 46 requisitos trazados, de los
    cuales el $100$\,\% de los 23 de prioridad \texttt{Must} cuenta con
    evidencia de verificación real (\autoref{cap:ingenieria-requisitos}) --
    sostenida, no solo declarada: la matriz de trazabilidad
    (\autoref{sec:matriz-resumen}) asigna un mecanismo de acceso a datos
    concreto a cada uno de los 46.
  \item Evidencia empírica reproducible y versionada como código: 5
    corridas de k6 con análisis estadístico formal (Wilcoxon, Cliff's
    delta), cobertura JaCoCo real ($82{,}97$\,\% de líneas, commit
    \texttt{0d1474d}), y auditoría de accesibilidad/rendimiento con
    Lighthouse -- las tres con datos crudos versionados en
    \texttt{docs/mediciones/}, no capturas de pantalla editadas a mano
    (\autoref{cap:resultados}).
  \item Un checklist de calidad de requisitos INCOSE~v4 aplicado con
    honestidad metodológica completa, documentando explícitamente qué
    requisitos no cumplen cada característica y por qué
    (\texttt{docs/checklists/incose2023-req.md}).
  \item Documentación arquitectónica versionada como código fuente: el
    modelo C4 completo en Structurizr DSL y 13 Architecture Decision
    Records (\autoref{cap:diseno-arquitectura}).
  \item Una revisión de trabajos relacionados con 30 referencias
    verificadas individualmente contra Crossref -- proceso que detectó y
    corrigió 5 errores de atribución de fuente antes de publicarse
    (\autoref{cap:trabajos-relacionados}) -- y una brecha de literatura
    identificada y sostenida con los resultados reales de este documento
    (\autoref{sec:disc-trabajos-relacionados}).
\end{itemize}

\section{Estado del proyecto al cierre de esta entrega}
\label{sec:concl-estado}

SGB-SaaS cierra esta Entrega Final como el software etiquetado
\texttt{v1.0.0} (portada de este documento), producto de 8 ramas de
módulos de dominio mergeadas a \texttt{main} durante el ciclo de esta
entrega (configuración paramétrica, préstamos avanzado, notificaciones y
verificación, catálogo avanzado, panel de administración, seguridad de
transporte, asistente con IA generativa, e infraestructura de CI), sobre
una base de 46 requisitos especificados y trazados, con el $100$\,\% de
los requisitos \texttt{Must} verificados
(\autoref{cap:ingenieria-requisitos}) y $82{,}97$\,\% de cobertura de
línea en 247 pruebas automatizadas del backend
(\autoref{sec:res-cobertura}). De los cinco bloques de evaluación
empírica exigidos por la guía, cuatro están completos con umbral
cumplido (rendimiento, seguridad, cobertura, calidad web), y uno --
usabilidad -- no se ha ejecutado todavía
(\autoref{tab:res-resumen}). El \autoref{cap:trabajo-futuro} deja
protocolizadas las cinco líneas concretas que cerrarían esos gaps, incluida
una vulnerabilidad de auto-degradación del rol \texttt{ADMIN} detectada
durante la propia redacción de este documento. Esta distribución -- logros
reales junto a gaps declarados sin ocultarlos -- es, en última instancia,
la conclusión más importante de este documento: no que SGB-SaaS esté
completamente terminado, sino que su estado real, con evidencia
verificable en cada afirmación, es conocido con precisión al cierre de
esta entrega.
